\chapter{Что такое искусственный интеллект}
\label{ch:intro}

В этой главе мы дадим неформальное представление о том, какие задачи
решает искусственный интеллект, и наметим математический «инструментарий»,
без которого современный ИИ не существовал бы.

\section{Краткая история идеи}

Современное представление об искусственном интеллекте складывалось
постепенно. В 1950~году Алан Тьюринг в статье \emph{Computing Machinery
and Intelligence} предложил операциональный критерий «мыслящей» машины
(тест Тьюринга), а сам термин \emph{искусственный интеллект} закрепился
после Дартмутского семинара 1956~года. Затем идею развивали несколько
сменявших друг друга подходов: \emph{экспертные системы} на правилах
(1970--1980-е годы), \emph{статистическое машинное обучение} на данных
(1990--2000-е), \emph{глубокое обучение} на нейронных сетях (с начала
2010-х) и, наконец, \emph{большие языковые модели} на основе архитектуры
трансформера (с 2017~года), определяющие облик ИИ сегодня.

\section{Чему мы научимся в этой книге}

В книге читателю встретятся три «слоя» материала:
\begin{enumerate}
  \item \textbf{Математические основы}: линейная алгебра, элементы
        анализа, теория вероятностей и численные методы оптимизации.
  \item \textbf{Алгоритмы машинного обучения}: от линейной регрессии
        до многослойных нейронных сетей.
  \item \textbf{Современный ИИ}: трансформеры, большие языковые модели,
        генеративные методы.
\end{enumerate}

В главе~\ref{ch:opt} мы подробно изучим один из старейших и до сих пор
важнейших численных методов — \emph{метод касательных Ньютона}, который
используется и для решения уравнений, и для обучения современных
нейросетей.
